%% ── Section 4: Experimental Setup ───────────────────────────
\section{Experimental Setup}
\label{sec:experiments}

\subsection{Dataset and Preprocessing}

We evaluate all frameworks on the \textbf{DDR benchmark dataset}~\cite{li2019ddr},
comprising 12,522 fundus photographs. Excluding ungradable Grade~5 images leaves
6,260 training, 2,503 validation, and 2,514 test images. Images are centre-cropped
to remove black borders and resized to $224 \times 224$ pixels. During local training,
standard augmentations (random horizontal flips, slight rotations) are applied to
prevent overfitting. The training set displays heavy natural class imbalance:
Grade~0 (50.0\%), Grade~1 (4.9\%), Grade~2 (36.0\%), Grade~3 (2.3\%), and Grade~4 (6.8\%).

\subsection{Modeling Clinical Heterogeneity}

We instantiate $K = 5$ clients under three statistical scenarios (visualised in
Fig.~\ref{fig:distributions}):
\begin{itemize}[leftmargin=*,noitemsep,topsep=2pt]
  \item \textbf{IID Partition:} Equal stratified random assignment (1,252 samples
  per client). This serves as the homogeneous theoretical upper bound for FL.
  \item \textbf{Dirichlet Non-IID ($\alpha=0.1$):} Pathological skew generated
  via a Dirichlet distribution $p_k \sim \mathrm{Dir}(\alpha)$. An $\alpha=0.1$
  forces each client to draw samples almost exclusively from a single class.
  Consequently, client sizes range from 32 to 2,087 samples; Client~0 contains
  solely 32 images (100\% Grade~0), while Client~4 holds 2,087 images dominated
  (98\%) by Grade~2.
  \item \textbf{Real-Site Clinical Partition:} We manually distribute data to model
  five realistic clinical archetypes: Site~0 (general screening clinic, G0-heavy,
  40\% volume), Site~1 (rural primary care, G1/G2-heavy, 22\%), Site~2 (specialist
  referral eye hospital, G3/G4-heavy, 10\%), Site~3 (diabetic retinopathy clinic,
  G2/G3-heavy, 19\%), and Site~4 (balanced urban hospital, 9\%).
\end{itemize}

\subsection{Network Architecture and Baselines}

All federated experiments utilise an EfficientNet-B3 backbone, chosen for its
optimal balance between parameter efficiency and feature extraction capability for
$224 \times 224$ resolution images. To simulate real-world deployment where
pre-trained weights are adapted, all FL runs warm-start from a centralised model
trained for 30 epochs (peak QWK = 0.8494). We benchmark OrdinalFed against:
(1)~\textbf{Centralised Baseline}, (2)~\textbf{FedAvg}~\cite{mcmahan2017fedavg},
and (3)~\textbf{FedProx}~\cite{li2020fedprox} with proximal coefficient $\mu = 0.01$.
All methods execute $T = 10$ communication rounds with $E = 5$ local epochs,
batch size 64, Adam optimiser ($\eta = 10^{-4}$), and use the same empirical
class-imbalance weights $\mathbf{c}$.
